% Generated draft: requires mathematical and source review.
\documentclass[11pt]{article}
\usepackage[margin=25mm]{geometry}
\usepackage{fontspec}
\setmainfont{Latin Modern Roman}
\usepackage{amsmath,amssymb,mathtools,unicode-math}
\setmathfont{Latin Modern Math}
\usepackage{enumitem,needspace,xurl,hyperref}
\setlength{\parindent}{0pt}
\setlength{\parskip}{5pt}
\setlength{\emergencystretch}{3em}
\newcommand{\N}{\mathbb{N}}
\newcommand{\Z}{\mathbb{Z}}
\newcommand{\Q}{\mathbb{Q}}
\newcommand{\R}{\mathbb{R}}
\newcommand{\C}{\mathbb{C}}
\newcommand{\F}{\mathbb{F}}
\newcommand{\A}{\mathbb{A}}
\newcommand{\PP}{\mathbb{P}}
\newcommand{\E}{\mathbb{E}}
\newcommand{\eps}{\varepsilon}
\newcommand{\dd}{\,\mathrm{d}}
\begin{document}
\section*{UnsolvedMath ID: 6800010}
\textbf{Source problem code:} catalogue entry
\subsection{Definitions and mathematical statement}
Let \(M\) be a connected closed smooth manifold, and let \(p:\widetilde M\to M\) be its universal covering. Assume \(\pi_1(M)\) is finite of odd order. Consequently the universal covering has finite odd degree. A metric has positive scalar curvature when its scalar curvature is strictly positive at every point.

Does the implication
\[
 \bigl[\exists\text{ a Riemannian metric }\widetilde g\text{ on }\widetilde M:
 \operatorname{Scal}_{\widetilde g}>0\bigr]
 \quad\Longrightarrow\quad
 \bigl[\exists\text{ a Riemannian metric }g\text{ on }M:
 \operatorname{Scal}_g>0\bigr]
\]
hold for every such \(M\)? The given positive-scalar-curvature metric on the cover need not be invariant under deck transformations; no invariance assumption is inserted. The statement asks for existence of some metric on the base, not descent of the particular metric supplied on the cover. No spin, dimension-at-least-five, toral, or other restriction is added from the section heading or related theory.

The supplied summary describes this as an unresolved odd-order-cover descent question and mentions Hanke/Stolz index-theoretic methods and \(\alpha\)-invariant descent as related tools. Its general warning that positive scalar curvature need not descend under arbitrary finite quotients is not a counterexample to the exact odd-order fundamental-group assertion. No complete resolution or precise supporting theorem is supplied; the open label remains a dated source assessment without independent current verification.
\subsection{Short English statement}
If a connected closed manifold has finite fundamental group of odd order and its universal cover supports positive scalar curvature, must the manifold itself support such a metric? The cover metric is not assumed invariant under the deck group.
\subsection{Sources}
\begin{itemize}
\item Local registry: \url{data/unsolvedmath/problems.json}, permanent UnsolvedMath ID \(6800010\), source code catalogue entry; supplied statement, background, and research summary, updated 2026-08-06. Proposed by Bernhard Hanke, Augsburg, 2018.
\item Morgan and Pansu, \emph{A List of Open Problems in Differential Geometry} (2018), Question environment 10, ``Toral manifolds and positive scalar curvature'': \url{https://www.imo.universite-paris-saclay.fr/~pansu/problems_MTDG.tex}. Supplied source-TeX extraction accessed 2026-07-29; current status marked as needing review.
\item Hanke, Stolz, Dirac-operator methods and \(\alpha\)-invariant descent are mentioned in the summary without precise original-paper citations or a resolving theorem.
\end{itemize}
\end{document}
